Before we bend space and curve time we must first understand why anyone would ever have thought to try. \textbf{The history of general relativity is not the history of a lone genius working in isolation, it is the history of nearly twenty five centuries of argument about a single, deceptively simple question}: why do bodies move the way they do? Every civilisation that has looked up at the night sky or dropped a stone from a height has had to answer this question in its own way, and the theory we are about to build in this book is, in a very real sense, the latest instalment of that argument rather than its final word. In this chapter we are going to walk through that argument from the Greeks to the physics of today, and along the way we will meet the geometers whose abstract mathematics quietly waited, sometimes for decades, for a physicist to notice that it was exactly what nature needed.

\section{Motion in the Ancient World}

    \subsection{The Aristotelian Cosmos}

    For the better part of two thousand years the Western understanding of motion belonged to Aristotle (384 to 322 B.C.), whose \textit{Physics} divided all motion into two families. \textit{Natural motion} was the motion a body underwent when left alone, and it depended entirely on what the body was made of: earth and water naturally fell towards the centre of the cosmos, while air and fire naturally rose away from it, and the heavenly bodies, made of a fifth incorruptible element he called \textit{aether}, naturally moved in perfect eternal circles. \textit{Violent motion}, by contrast, was any motion imposed on a body against its nature, and Aristotle's central postulate about it is one every student of mechanics eventually collides with: \textbf{a body moves only while something is in contact with it and pushing it}, so that the moment the pushing stops the motion must stop too \citep{aristotle_physics}. This is why a thrown spear troubled him so much. If the hand that threw it is no longer touching it, what keeps it moving through the air after it leaves the fingers? Aristotle's answer, the doctrine of \textit{antiperistasis}, was that the air parted by the spear's tip rushes around to its rear and pushes it onwards, a mechanism that is elegant on paper and hopeless the moment one actually thinks about throwing something, and it would take the better part of two millennia and the combined efforts of Alexandria, Baghdad and Paris to finally dismantle it.

    Underneath this physics of pushes and pulls lay a much grander architecture, the geocentric cosmos, in which a stationary Earth sat at the centre of a finite universe and everything else, the Moon, the Sun, the planets and the fixed stars, was carried around it on nested, concentric, crystalline spheres. It is worth pausing on how seriously this was meant: for Aristotle the centre of the cosmos was not merely a convenient coordinate origin, it was the natural resting place towards which heavy matter was drawn, so that geometry and physics were one and the same fact.

    \subsection{Ptolemy and the Mathematics of the Heavens}

    The trouble with perfect circles around a stationary Earth is that the planets simply refuse to move on them, drifting faster and slower across the year and occasionally, most embarrassingly of all, reversing direction in the sky entirely, a phenomenon we now call retrograde motion and understand as nothing more than the view from a moving Earth overtaking a slower outer planet. Claudius Ptolemy, working in Alexandria around 150 A.D., inherited several centuries of Greek attempts to rescue the circles (Apollonius and Hipparchus among them) and assembled them into the \textit{Almagest}, the most complete and mathematically serious astronomical text the ancient world produced \citep{ptolemy_almagest}. Ptolemy kept the philosophical commitment to circular motion but smuggled in three devices to make the mathematics fit the sky: the \textit{epicycle}, a small circle whose centre rides upon a larger circle called the \textit{deferent}; the \textit{eccentric}, which quietly moves the Earth slightly off the true centre of the deferent; and, most cunning of all, the \textit{equant}, an empty point in space from which the planet's motion \textit{appears} uniform even though it is uniform nowhere else. The equant in particular is a fascinating piece of intellectual dishonesty dressed up as rigour, because it abandons the very principle of uniform circular motion it was invented to save, and it is precisely this point, as we shall see in a moment, that the astronomers of Islam would spend centuries trying to repair.

\section{The Arabic Custodians of Motion}

    When the intellectual centre of gravity of the ancient world shifted from Alexandria towards Baghdad, Damascus and later Maragha and Samarkand, it carried Aristotle and Ptolemy with it, but it did not simply preserve them, it interrogated them. The crucial bridge figure is not, strictly speaking, an Arab at all: John Philoponus, a Christian philosopher working in sixth century Alexandria, wrote a commentary on Aristotle's \textit{Physics} in which he attacked antiperistasis directly, proposing instead that the thrower impresses onto the projectile itself a finite, internal quantity of motive power, an \textit{incorporeal kinetic force}, which then dissipates gradually as the projectile travels, explaining why a spear eventually falls without needing the air to conspire behind it. Philoponus's Greek text was largely forgotten in Byzantium but survived and was studied intensely by scholars writing in Arabic, and it is through that transmission, not through the Latin West, that his idea re-entered the mainstream of physics.

    The Persian polymath Ibn Sina, known in Latin as Avicenna (980 to 1037), sharpened Philoponus's idea considerably in his encyclopaedic \textit{Kitab al-Shifa} (The Book of Healing). He argued that the impressed force, which he called \textit{mayl} (inclination), is a genuine, quantifiable property of the moving body itself rather than of the medium around it, and, in a passage that reads almost eerily like a draft of Newton's first law, he claimed that in a void this \textit{mayl} would never be consumed and the body would continue moving forever, since nothing in a true vacuum could exist to wear it down \citep{rashed1996}. This is not yet inertia in the modern sense, because Ibn Sina still thought of \textit{mayl} as a cause that must be continuously present to sustain motion rather than as the mere absence of a cause to stop it, but the intuition that unresisted motion is self-perpetuating had, for the first time, been stated clearly and defended on physical grounds rather than dismissed as an impossible thought experiment about an empty universe Aristotle himself had declared could not exist.

    Astronomy saw an even more consequential development. Frustrated by the equant's abandonment of true uniform circular motion, the mathematicians of the Maragha observatory in thirteenth century Persia, led by Nasir al-Din al-Tusi, invented a purely geometrical device, now called the \textit{Tusi couple}, a small circle rolling inside a larger one, whose combined motion can reproduce a straight line oscillation using only uniformly rotating circles, with no equant required. A century later in Damascus, Ibn al-Shatir used exactly this device to rebuild the entire Ptolemaic planetary system without a single equant, producing models that were mathematically superior to Ptolemy's own \citep{saliba1987}. What makes this genuinely remarkable, and worth pausing on for a student of physics rather than only of history, is that when Nicolaus Copernicus published \textit{De Revolutionibus Orbium Coelestium} in 1543, the geometrical diagrams he used to eliminate the equant from his own, now heliocentric, system are mathematically identical to Ibn al-Shatir's, down to the labelling of points in some manuscripts. Historians still debate exactly how the Maragha school's mathematics reached Copernicus, whether through Byzantine Greek intermediaries in Italy or some now lost channel, but the technical debt is not seriously disputed, and it is a useful reminder that the Copernican revolution did not spring from nowhere: it stood, quite literally, on Islamic geometrical shoulders.

\section{The Latin West Breaks with Aristotle}

    \subsection{Buridan's Impetus}

    The idea of an impressed motive force crossed into Latin Europe through translations of Arabic and Greek texts flooding into the universities of the thirteenth century, and it found its sharpest medieval expression in Jean Buridan, rector of the University of Paris in the 1340s. Buridan proposed that the mover impresses upon the projectile a genuine, permanent quantity he called \textit{impetus}, proportional to the amount of matter in the body and to its speed (a description that should already sound suspiciously familiar to anyone who has met the word \textit{momentum}), and that this impetus would in principle persist forever were it not continuously eroded by air resistance and by the body's own natural tendency to fall \citep{clagett1959}. Buridan went further than Ibn Sina in one crucial respect: he applied impetus theory to the heavens themselves, suggesting that God, at the Creation, might simply have impressed upon the celestial spheres an impetus that meets no resistance in the frictionless heavens and therefore never decays, a proposal that quietly detaches the eternal circular motion of the planets from any need for a continuously acting mover, Aristotle's prime mover included. His student Nicole Oresme pushed the mathematics further still, inventing a graphical method of plotting a body's speed against time as a geometrical figure and showing that the area under such a figure gives the distance travelled, a genuinely startling anticipation, two and a half centuries early, of the mean speed theorem that Galileo would later use to characterise uniformly accelerated fall.

    \subsection{Copernicus and Kepler}

    Copernicus's 1543 relocation of the Sun to the centre of the cosmos was, in itself, primarily a mathematical and aesthetic argument (he found the resulting system simpler and more harmonious), and it would take Johannes Kepler, working with the extraordinarily precise naked eye observations that Tycho Brahe had spent a lifetime accumulating, to turn heliocentrism into a genuinely predictive physical theory. Between 1609 and 1619 Kepler announced three laws that any general relativity student should keep close at hand, because they are the observational bedrock that Newton, and much later Einstein, would both have to reproduce: planets move on ellipses with the Sun at one focus, they sweep out equal areas in equal times, and the square of their orbital period is proportional to the cube of their orbit's semi-major axis. Kepler arrived at the ellipse only after years of fighting his own preference for circles, famously calling his discovery a \textit{cartload of dung} in a fit of frustration with the messy, imperfect shape nature had apparently chosen over the geometrical perfection he had hoped for, and it is a nice historical irony that this same ellipse would reappear, itself only an approximation, when Einstein finally explained in 1915 why Mercury's own elliptical orbit slowly, stubbornly, refuses to close.

    \subsection{Galileo and the First Principle of Relativity}

    It is with Galileo Galilei that impetus theory finally becomes something closer to modern inertia. In his \textit{Two New Sciences} of 1638, written under house arrest and smuggled out of Italy for publication in the Protestant Netherlands, Galileo established through careful (and, for the tools available to him, remarkably precise) experiments with balls rolling down inclined planes that a freely falling body accelerates uniformly, covering distances in the ratio of the odd numbers one, three, five and so on in successive equal intervals of time, and he reasoned that a body moving on a perfectly smooth horizontal surface, meeting no resistance and no incline either to speed it up or slow it down, would continue moving forever at constant speed \citep{galileo1638}. This is, for all practical purposes, Newton's first law stated half a century early.

    Just as important for our story, though, is an idea Galileo had already set out six years earlier, in 1632, in the \textit{Dialogue Concerning the Two Chief World Systems}. There he asks the reader to imagine themselves below decks on a smoothly sailing ship, with no portholes, watching fish swim in a bowl, butterflies fly and drops of water fall from a bottle, and he argues that not one of these phenomena would betray to the observer whether the ship is standing still in harbour or sailing at constant velocity across a perfectly calm sea \citep{galileo1632}. \textbf{No mechanical experiment performed entirely inside a closed cabin can distinguish uniform motion from rest}, and this is the first clean statement of what we now call the \textit{principle of relativity}, the very principle that, generalised first to electromagnetism by Einstein in 1905 and then to accelerated frames by the equivalence principle we shall meet in chapter 3, becomes the seed from which the entire theory of general relativity eventually grows. Tradition, almost certainly apocryphal but too delicious to leave out, has Galileo muttering \textit{E pur si muove} (``and yet it moves'') under his breath after being forced by the Roman Inquisition in 1633 to formally recant the Copernican system his own physics had done so much to defend.

\section{Newton and the Absolute Stage}

    Isaac Newton's \textit{Philosophiae Naturalis Principia Mathematica}, published in 1687, is the book that finally welded Galileo's terrestrial kinematics to Kepler's celestial geometry into a single mechanics. Newton's three laws of motion, and above all his law of universal gravitation, showed that the very same force that pulls an apple to the ground also holds the Moon in its orbit and drives the planets along Kepler's ellipses, a genuinely staggering unification for its time \citep{newton1999principia}. But Newton needed a stage on which this mechanics could play out, and he built one deliberately: \textit{absolute space}, an unmoving, featureless, infinite container existing independently of anything inside it, and \textit{absolute time}, flowing uniformly and identically for every observer everywhere in the universe. He even offered an experimental argument for taking absolute space seriously, the famous rotating bucket: spin a bucket of water and its surface climbs the sides into a concave shape, and since this happens whether or not the bucket is rotating relative to any other nearby matter, Newton argued that the water must be rotating relative to something, namely space itself.

    What is often forgotten by students who only encounter Newton through the tidy inverse square law is how deeply uncomfortable Newton himself was with one consequence of his own theory: gravity, as he had written it down, acts instantaneously at a distance across empty space, with no mediating mechanism whatsoever, and this troubled him far more than it troubles most introductory textbooks today. In a 1693 letter to the theologian Richard Bentley he wrote, with striking force, that

    \begin{quote}
    ``that gravity should be innate, inherent and essential to matter, so that one body may act upon another at a distance through a vacuum, without the mediation of any thing else, by and through which their action and force may be conveyed from one to another, is to me so great an absurdity that I believe no man who has in philosophical matters a competent faculty of thinking can ever fall into it.'' \citep{newton1958papers}
    \end{quote}

    Newton had no candidate mechanism to offer in its place, and so, with characteristic caution, he famously declined to speculate publicly, writing in the General Scholium he added to later editions of the \textit{Principia} that \textit{hypotheses non fingo}, ``I feign no hypotheses'' \citep{newton1999principia}. It would take more than two centuries, and an entirely new branch of geometry that did not yet exist in Newton's lifetime, before anyone could offer a mechanism at all: general relativity's answer, that gravity is not a force transmitted through space but the curvature of space-time itself, is in a very real sense the belated reply to Newton's own unease.

\section{The Geometers Who Prepared the Ground}

    \subsection{Gauss and the Curvature Hidden Inside a Surface}

    For gravity to become geometry, geometry first had to learn how to describe curvature without appealing to anything outside itself, and this is exactly what Carl Friedrich Gauss achieved in 1827 in his \textit{Disquisitiones Generales circa Superficies Curvas} (General Investigations of Curved Surfaces) \citep{gauss1828}. Gauss showed that the curvature of a surface, a quantity that seems at first to depend entirely on how that surface bends through the three dimensional space surrounding it, can in fact be computed using only measurements made \textit{within} the surface itself, lengths and angles as measured by an ant living on it, with no reference whatsoever to the ambient space. He was so struck by this fact that he named it his \textit{Theorema Egregium}, the ``remarkable theorem'', and remarkable it is: it is the very first proof that curvature is an \textit{intrinsic} property of a space, not something a space merely borrows from a higher dimensional stage on which it happens to sit. This is precisely the property that will let us, in chapter 2, speak of the curvature of space-time without ever having to ask what space-time is supposedly curving \textit{into}.

    It is often repeated, and worth mentioning with the appropriate degree of caution, that Gauss was not only a mathematician but also, in his day job, the director of a two decade long geodetic survey of the Kingdom of Hanover, and that he measured the angles of an enormous triangle formed by the peaks of the Brocken, the Hohenhagen and the Inselsberg, partly to test whether the angles of a very large physical triangle really do sum to exactly 180 degrees, as flat Euclidean geometry demands, or whether the physical space we live in might, just possibly, be very slightly curved. Whether Gauss genuinely intended this as a test of the geometry of physical space, or whether the story has been embellished by later admirers eager to cast him as a prophet of general relativity, remains disputed among historians, and no deviation from 180 degrees was in any case ever detected, the curvature of space-time being far too gentle at the scale of a mountain range for eighteenth century theodolites to have any hope of catching it. What is beyond dispute is that Gauss had, at the very least, begun to entertain the question of whether Euclidean geometry is a truth of pure reason or an empirical fact about our particular universe, and that question is the hinge on which this entire chapter, and indeed this entire book, turns.

    \subsection{Riemann's Prophecy}

    Gauss's intrinsic geometry described curved surfaces, objects with only two dimensions. The decisive generalisation, to spaces of any number of dimensions whatsoever, was the achievement of Gauss's own student, Bernhard Riemann, and it arrived in a single lecture. In 1854, required by University of Göttingen custom to deliver a probationary lecture (a \textit{Habilitationsvortrag}) before the assembled faculty in order to earn the right to teach, the shy and sickly Riemann chose, on Gauss's own suggestion, the topic ``\"{U}ber die Hypothesen, welche der Geometrie zu Grunde liegen'' (On the Hypotheses which lie at the Bases of Geometry). Gauss himself, by then an old man near the end of his life and famously difficult to impress, sat in the audience, and by every surviving account he was profoundly moved by what he heard.

    In that lecture Riemann introduced what we would now call an $n$-dimensional manifold, a space that need only look locally like ordinary flat Euclidean space, and he showed how to equip such a space with a quantity, the \textit{metric}, encoding the distance between any two infinitesimally close points, from which a notion of curvature at every point could then be built by pure calculation, generalising Gauss's Theorema Egregium to any dimension and to spaces nobody could ever hope to visualise by embedding them in something larger. This is the mathematical object we will meet formally as $g_{\mu\nu}$ in chapter 2, and without it there simply is no general relativity to write down. But Riemann did something else that evening, something far stranger for a lecture ostensibly about pure mathematics, something translated from the German by the English mathematician William Kingdon Clifford two decades later and read ever since by physicists with a peculiar chill of recognition: he wondered aloud, in the lecture's closing paragraphs, whether the geometry of physical space might not be fixed and given after all, but might instead be determined by the matter distributed within it. \begin{quote}
    ``we must seek the ground of its metric relations outside it, in binding forces which act upon it'' \citep{riemann1873}
    \end{quote} is how Clifford rendered Riemann's own conjecture, and Riemann explicitly named this a question for physics, not mathematics, to eventually answer, adding that the answer could only come from experience and observation of the actual physical world. Riemann died of tuberculosis in 1866, at only thirty nine years old, on the shores of Lake Maggiore in Italy, sixty years before anyone would show that he had, in that single lecture, more or less written down the mathematical vessel into which Albert Einstein would one day pour a physical theory of gravity.

\section{Einstein's Long Road to General Relativity}

    Albert Einstein's 1905 theory of special relativity had already overthrown Newton's absolute time, welding space and time together into a single four dimensional structure and showing that simultaneity itself depends on the observer's motion. But special relativity, as its name honestly admits, is special: it applies only to unaccelerated, inertial observers, and it could offer no account whatsoever of gravity, since Newtonian gravity's instantaneous action at a distance is flatly incompatible with special relativity's insistence that no signal, gravitational or otherwise, can outrun the speed of light.

    The idea that would eventually resolve this came to Einstein in 1907, while he was still working at the Swiss patent office in Bern, and he later called it, in a 1922 lecture in Kyoto, \textit{der gl\"{u}cklichste Gedanke meines Lebens}, ``the happiest thought of my life'': a person falling freely from a rooftop feels no gravity whatsoever, their own weight vanishing entirely for as long as the fall lasts. From this simple, almost childlike observation Einstein extracted his \textit{equivalence principle}, the proposal that being at rest in a gravitational field is locally indistinguishable from accelerating through empty space with no gravity at all, a principle that, as we shall see explicitly in chapter 3, is nothing less than Galileo's shipboard relativity extended from uniform motion to acceleration itself, and it is the true conceptual seed of everything that follows in this book.

    Turning that seed into a working theory took Einstein eight more gruelling years. He understood by around 1912, while teaching in Prague, that a theory built on the equivalence principle would need a gravitational field capable of bending the paths of light rays and of clocks running at different rates in different places, and that flat Minkowski space-time (the space-time of special relativity) simply could not accommodate this. What he needed was some kind of curved geometry, and he had no idea such a thing already existed in full generality, sitting in the mathematics literature since 1854, until his old university classmate from Zurich, the mathematician Marcel Grossmann, came to his rescue. According to Einstein's own later recollection, and as reported by his biographer Abraham Pais, Einstein turned to Grossmann in something close to desperation with the words \textit{``Grossmann, du musst mir helfen, sonst werd' ich verr\"{u}ckt!''} (``Grossmann, you must help me or I shall go crazy!'') \citep{pais1982}. Grossmann answered by introducing Einstein to precisely Riemann's geometry, together with the \textit{absolute differential calculus} that the Italian mathematicians Gregorio Ricci-Curbastro and Tullio Levi-Civita had spent the intervening decades building on top of it, the very tensor calculus this book will spend the whole of chapter 2 constructing from scratch.

    In 1913 Einstein and Grossmann published together a first attempt, the so called \textit{Entwurf} (draft) theory, which used Riemann's curved metric but, out of a mixture of genuine physical worry and outright mathematical error, stopped short of demanding that the field equations hold in every coordinate system \citep{einsteingrossmann1913}. The Entwurf theory predicted the wrong value for the century old, already precisely measured anomaly in the precession of Mercury's perihelion, and Einstein knew it. The final months, in Berlin during the autumn of 1915, were by his own account among the most intense of his life: in a sequence of four papers delivered to the Prussian Academy of Sciences on four successive Thursdays in November, he corrected his errors in full public view, week by week, until on 25 November 1915 he presented the field equations essentially in the form we use today \citep{einstein1915}, and applied them almost immediately to recover Mercury's precession exactly, a moment he later described as having given him heart palpitations severe enough to worry he might be physically unwell.

\section{Riemann, Hilbert and the Question of Theft}

    It is sometimes claimed, in the looser corners of popular science writing, that Einstein ``stole'' general relativity from Riemann, and it is worth being precise about why this particular accusation does not survive contact with the historical record, even though a genuine and fascinating priority dispute does sit at the heart of this story, just with a different antagonist. Riemann died in 1866, forty nine years before Einstein wrote down a single field equation, and he never worked on gravity, never proposed a field equation, and never claimed to have solved the problem this chapter has been tracing since Aristotle: Riemann supplied a mathematical \textit{language}, the geometry of curved manifolds, not a physical theory of gravitation, and Einstein always said so plainly and gratefully, both in his 1913 paper with Grossmann and later in his own \textit{Autobiographical Notes} \citep{einstein1949}. Crediting Riemann with the mathematics he actually built, generously and by name, is simply accurate history, not theft.

    The real, and rather more dramatic, priority dispute concerns not Riemann but David Hilbert, the towering G\"{o}ttingen mathematician who was, quite independently, hunting for the same field equations in the very same autumn of 1915. Einstein had visited G\"{o}ttingen that summer and lectured on his still unfinished theory, and Hilbert, characteristically supremely confident (he is reported to have remarked, only half in jest, that \textit{``every boy in the streets of G\"{o}ttingen understands more about four dimensional geometry than Einstein''}, so thoroughly had David Hilbert's own city absorbed the mathematics Einstein was still struggling to apply), set to work on the problem from a different direction entirely, a variational principle built from the very Einstein-Hilbert action we shall derive in chapter 3. Hilbert submitted a paper to the G\"{o}ttingen Academy titled ``Die Grundlagen der Physik'' on 20 November 1915, five days before Einstein's own final, decisive paper of 25 November \citep{hilbert1915}, and for decades afterwards a persistent rumour circulated that Einstein had somehow seen and silently borrowed from Hilbert's unpublished proofs, or conversely that Hilbert had quietly lifted the final covariant equations from Einstein's own G\"{o}ttingen lecture.

    The matter was only properly settled in 1997, when the historians of science Leo Corry, J\"{u}rgen Renn and John Stachel tracked down the actual surviving galley proofs of Hilbert's paper, printer's proofs dated December 1915, sitting quietly in the G\"{o}ttingen archive \citep{corry1997}. Those proofs, it turned out, did \textit{not} yet contain the fully generally covariant field equations that appeared in Hilbert's paper when it was finally published, months later, in 1916: Hilbert had revised his own manuscript after that publication date, almost certainly after seeing Einstein's completed November papers, to bring it into line with Einstein's result. \textbf{Einstein reached the correct, generally covariant field equations first, and independently}, while Hilbert arrived at essentially the same equations by a beautiful and genuinely independent variational route, close behind him and very possibly influenced, in the final polish of his own paper, by Einstein's success. Hilbert himself, to his considerable credit, never once publicly disputed Einstein's priority and always referred to the result as ``Einstein's equations'', and the two men remained on entirely cordial, mutually respectful terms for the rest of their lives. If there is a lesson here for a student just beginning this subject, it is that the discovery of general relativity was not a single flash of individual genius but the convergence of an entire community, Gauss's intrinsic curvature, Riemann's manifolds, Ricci and Levi-Civita's calculus, Grossmann's tutoring, and Hilbert's own parallel and generous rivalry, all funnelled through Einstein's singular physical insight that gravity simply \textit{is} the shape of space-time.

\section{Confirmation: From the Trenches to the Eclipse}

    Extraordinary claims of this kind demand to be tested, and their testing produced two of the most romantic episodes in the whole history of physics. The first came almost immediately. Karl Schwarzschild, a distinguished German astrophysicist who had volunteered for military service and was at the time calculating artillery trajectories on the Russian front, read Einstein's November 1915 papers and, within a matter of weeks, working between his military duties, found the first exact, non approximate solution to Einstein's notoriously difficult field equations, the solution describing the curved space-time around a single spherical, non rotating mass, which we will study closely in chapter 3 and which today bears his name \citep{schwarzschild1916}. Schwarzschild mailed his solution to Einstein directly from the front; Einstein, visibly delighted, presented it to the Prussian Academy on Schwarzschild's behalf in January 1916. Schwarzschild had, tragically, already contracted a rare and painful autoimmune skin disease at the front, and he died only a few months later, in May 1916, never learning that his solution would one day predict the existence of black holes.

    The second episode made Einstein a global celebrity almost overnight. General relativity predicted that starlight passing close to the Sun should be deflected by the Sun's gravity by twice the amount Newtonian gravity alone would predict, an effect only observable during a total solar eclipse, when the Sun's own overwhelming brightness is briefly blotted out and stars that ordinarily hide directly behind it become visible. On 29 May 1919, two expeditions organised chiefly by the British astronomers Arthur Eddington and Frank Dyson, one to the island of Pr\'{i}ncipe off the coast of West Africa and one to Sobral in northern Brazil, photographed the stars around the eclipsed Sun and found, after months of careful measurement, a deflection matching Einstein's relativistic prediction and decisively ruling out the smaller Newtonian value \citep{dyson1920}. When the results were announced at a joint meeting of the Royal Society and the Royal Astronomical Society in London that November, the news spread with astonishing speed: the \textit{New York Times} greeted it on 10 November 1919 with the headline ``Lights All Askew in the Heavens'', and within weeks Einstein, until then a name known chiefly to physicists, had become the most famous scientist on Earth.

\section{Beyond General Relativity}

    The theory did not stop working once it was confirmed, and it has continued, ever since, to reveal consequences even its own author found hard to accept. In 1917 Einstein had inserted an extra term, the \textit{cosmological constant}, into his own field equations purely to force them to describe a static, unchanging universe, since a static cosmos was simply what everyone, Einstein included, assumed the real universe to be. The Russian mathematician Alexander Friedmann in 1922 \citep{friedmann1922} and, independently, the Belgian priest and astronomer Georges Lema\^{i}tre in 1927 \citep{lemaitre1927}, showed that Einstein's own unmodified equations naturally describe an expanding, evolving universe instead, a result Einstein initially, and famously, resisted. It took Edwin Hubble's 1929 observation that distant galaxies are indeed systematically receding from us, and at a rate proportional to their distance, exactly as Friedmann and Lema\^{i}tre had predicted, to settle the matter \citep{hubble1929}, and Einstein is widely (if perhaps not entirely reliably) reported to have called his own cosmological constant the biggest blunder of his career.

    General relativity's most extreme prediction took even longer to be taken seriously. In 1939 J. Robert Oppenheimer and his student Hartland Snyder showed mathematically that a sufficiently massive star, having exhausted its nuclear fuel, must collapse without limit under its own gravity, its surface eventually crossing a point of no return beyond which not even light can climb back out \citep{oppenheimersnyder1939}, but the paper was largely ignored for nearly three decades, overshadowed by the Second World War and by a lingering, widely shared conviction among physicists, Einstein among them, that nature would surely find some way to forbid so pathological an outcome. It was the physicist John Archibald Wheeler who, in a 1967 lecture, popularised a short, vivid name for these collapsed objects, \textit{black holes}, a term that stuck instantly where more cumbersome earlier phrases such as ``gravitationally completely collapsed object'' never had \citep{wheeler1998}, and it was Wheeler, too, who gave us the aphorism already quoted on this chapter's own title page: spacetime tells matter how to move, and matter tells spacetime how to curve.

    Einstein's 1916 prediction of gravitational waves, tiny ripples of curvature radiating outward from accelerating masses, waited even longer for its confirmation, precisely one century, until on 14 September 2015 the two detectors of the Laser Interferometer Gravitational Wave Observatory (LIGO) simultaneously registered the unmistakable signature of two black holes, thirty six and twenty nine times the mass of the Sun, spiralling into one another and merging over a billion years ago, a discovery announced to the world in February 2016 and one we shall return to properly in chapter 4 \citep{abbott2016}.

    And there the story remains genuinely, actively unfinished. General relativity describes gravity magnificently at the scale of planets, stars and galaxies, yet it still refuses to speak the same language as quantum mechanics, the theory that governs the very small, and no one yet possesses a fully working theory of \textit{quantum gravity} capable of describing, for instance, what truly happens at the very centre of a black hole or in the first unimaginably brief instant of the universe's own existence. String theory, loop quantum gravity and a handful of other, still more speculative approaches are today's competing attempts to write the next chapter, each borrowing, in one way or another, from the very same geometrical inheritance that runs through this one: Gauss's intrinsic curvature, Riemann's manifolds, and Einstein's insight that geometry itself can be physics. It is worth carrying this long procession of names in mind as we turn, in chapter 2, to build the mathematics properly from first principles, because every symbol we are about to write down, the metric, the connection, the curvature tensor, was earned by someone, often across a lifetime and occasionally at real personal cost, and general relativity is best understood not as a single flash of individual genius but as the latest, and very probably not the last, chapter in humanity's long and stubborn argument about why things fall.
